% =====================================================================
%  Chapter 6 — The Lovász Local Lemma
%  File: chapters/06_lll.tex
% =====================================================================
\providecommand{\E}{\mathbb{E}}
\providecommand{\Prob}{\mathbb{P}}
\providecommand{\whp}{\textit{whp}}

\chapter{The Lov\'asz Local Lemma}
\label{chap:lll}

We met the local lemma in Chapter~1 (Theorem~\ref{thm:lll}), where it improved
the Ramsey lower bound, and promised a proof here. The setting is the awkward
middle ground between two extremes. If the bad events $A_1,\dots,A_n$ have
$\sum_i\Prob(A_i)<1$, the union bound avoids them all; if they are mutually
independent, $\Prob(\bigcap_i\overline{A_i})=\prod_i(1-\Prob(A_i))>0$ does. The
local lemma handles the common case in between --- many events, each depending on
only a few others --- and concludes avoidance is possible even when the union
bound fails badly.

% =====================================================================
\section{Statement and proof}
% =====================================================================

\begin{definition}{Independence from a set; dependency graph}{dependency}
An event $A$ is \emph{independent from} a set of events $\mathcal B$ if
$A$ is independent of every event in the $\sigma$-algebra generated by
$\mathcal B$. A graph $D$ on vertex set $[n]$ is a \emph{dependency graph} for
$A_1,\dots,A_n$ if each $A_i$ is independent from
$\{A_j:j\notin\{i\}\cup N(i)\}$, where $N(i)$ is the neighbourhood of $i$.
\end{definition}

In the common \emph{variable model}, each $A_i$ is determined by a subset $B_i$
of independent underlying variables, and joining $i\sim j$ whenever
$B_i\cap B_j\ne\emptyset$ yields a valid dependency graph.

\begin{theorem}{Local lemma, symmetric form}{lll-sym}
Let $A_1,\dots,A_n$ be events with $\Prob(A_i)\le p$, each independent from all
but at most $d$ of the others. If $e\,p\,(d+1)\le1$, then with positive
probability none of the $A_i$ occur.
\end{theorem}

\begin{theorem}{Local lemma, general form}{lll-general}
With $N(i)$ as in Definition~\ref{def:dependency}, suppose there are
$x_1,\dots,x_n\in[0,1)$ with
\[
  \Prob(A_i)\le x_i\prod_{j\in N(i)}(1-x_j)\qquad\text{for all }i.
\]
Then $\displaystyle\Prob\Bigl(\bigcap_i\overline{A_i}\Bigr)
\ge\prod_{i=1}^n(1-x_i)>0.$
\end{theorem}

\begin{proof}
We show, by induction on $|S|$, that
$\Prob\bigl(A_i\mid\bigcap_{j\in S}\overline{A_j}\bigr)\le x_i$ for every
$i\notin S$. Split $S=S_1\cup S_2$ with $S_1=S\cap N(i)$. Writing
$B=\bigcap_{j\in S_1}\overline{A_j}$ and $C=\bigcap_{j\in S_2}\overline{A_j}$,
\[
  \Prob\bigl(A_i\mid B\cap C\bigr)
  =\frac{\Prob(A_i\cap B\mid C)}{\Prob(B\mid C)}.
\]
The numerator is at most $\Prob(A_i\mid C)=\Prob(A_i)\le
x_i\prod_{j\in N(i)}(1-x_j)$, since $A_i$ is independent from $C$. For the
denominator, applying the inductive hypothesis to the events indexed by $S_1$
(each conditioned on a subset of $\overline{A_j}$'s) gives
$\Prob(B\mid C)\ge\prod_{j\in S_1}(1-x_j)\ge\prod_{j\in N(i)}(1-x_j)$. Dividing
leaves $\Prob(A_i\mid B\cap C)\le x_i$. Finally,
$\Prob(\bigcap_i\overline{A_i})=\prod_{i}\bigl(1-\Prob(A_i\mid\bigcap_{j<i}
\overline{A_j})\bigr)\ge\prod_i(1-x_i)$.
\end{proof}

\begin{proof}[Symmetric form from the general form]
Set every $x_i=1/(d+1)<1$. Then
$x_i\prod_{j\in N(i)}(1-x_j)\ge\frac1{d+1}\bigl(1-\frac1{d+1}\bigr)^d
\ge\frac1{e(d+1)}\ge p\ge\Prob(A_i)$, using $(1-\frac1{d+1})^d\ge e^{-1}$ and
the hypothesis $ep(d+1)\le1$.
\end{proof}

A frequently convenient repackaging asks only that each neighbourhood carry
little total probability.

\begin{corollary}{Neighbourhood form}{lll-nbhd}
If $\Prob(A_i)<\tfrac12$ and $\sum_{j\in N(i)}\Prob(A_j)\le\tfrac14$ for all $i$,
then with positive probability none of the $A_i$ occur.
\end{corollary}

\begin{proof}
Apply Theorem~\ref{thm:lll-general} with $x_i=2\Prob(A_i)$; then
$x_i\prod_{j\in N(i)}(1-x_j)\ge2\Prob(A_i)\bigl(1-\sum_{j\in N(i)}x_j\bigr)
\ge2\Prob(A_i)\cdot\tfrac12=\Prob(A_i)$.
\end{proof}

\begin{insight}
The general form beats the symmetric one by letting the slack $x_i$ differ,
which matters when the events have very different probabilities or degrees. The
proof turns on one bound: conditioning on other events not occurring can raise
$\Prob(A_i)$, but only through $A_i$'s neighbours and only up to $x_i$. Shearer
showed the constant $e$ is optimal; in practice its value rarely matters.
\end{insight}

% =====================================================================
\section{Colouring hypergraphs}
% =====================================================================

Chapter~1 showed (Theorem~\ref{thm:hyper_lower}) that a $k$-uniform hypergraph
with fewer than $2^{k-1}$ edges is $2$-colourable --- a \emph{global} bound on the
number of edges. The local lemma replaces it with a \emph{local} condition on
how edges overlap, which is far weaker to satisfy.

\begin{theorem}{Local $2$-colourability}{hypergraph-lll}
A $k$-uniform hypergraph is $2$-colourable if every edge intersects at most
$e^{-1}2^{k-1}-1$ other edges.
\end{theorem}

\begin{proof}
Colour vertices independently and uniformly, and let $A_e$ be the event that
edge $e$ is monochromatic, so $\Prob(A_e)=2^{1-k}$. Since $A_e$ depends only on
the colours of $e$'s vertices, it is independent of all edges disjoint from
$e$; if $e$ meets at most $d$ others then the symmetric local lemma applies once
$e\cdot2^{1-k}(d+1)\le1$, i.e.\ $d+1\le e^{-1}2^{k-1}$. Then some colouring
leaves no edge monochromatic.
\end{proof}

\begin{corollary}{Regular hypergraphs}{kregular}
For $k\ge9$, every $k$-uniform $k$-regular hypergraph is $2$-colourable.
\end{corollary}

\begin{proof}
Each edge shares a vertex with at most $k(k-1)$ others (its $k$ vertices lie in
$k-1$ further edges each), and $k(k-1)+1\le e^{-1}2^{k-1}$ holds for $k\ge9$.
\end{proof}

\begin{insight}
The move is to trade a global budget for a local one.
Theorem~\ref{thm:hyper_lower} caps the total number of edges, so it says nothing
about a hypergraph with many but loosely-interacting edges. The local lemma
ignores the edge count entirely and asks only that no edge overlap too many
others. Sparse interaction, not few objects, is what it needs --- which is why it
turns up throughout colouring, scheduling, and constraint satisfaction.
\end{insight}

% =====================================================================
\section{Latin transversals}
% =====================================================================

An $n\times n$ array with entries (``symbols'') has a \emph{Latin transversal} if
one can choose $n$ cells, one per row and one per column, all bearing distinct
symbols. Whether every Latin square has one is a famous open problem; a robust
positive result follows from a \emph{lopsided} strengthening of the local lemma.

\begin{theorem}{Erd\H{o}s--Spencer 1991}{latin-transversal}
Every $n\times n$ array in which each symbol appears at most $n/(4e)$ times has a
Latin transversal.
\end{theorem}

\begin{insight}
This was the first use of the lopsided local lemma. The conflicts --- two chosen
cells sharing a symbol --- are negatively associated: forcing one coincidence
makes the others less likely, since they compete for the same permutation.
Noticing that one-sided dependence is enough is the whole point; the lemma is
really about the direction of influence between events, not strict
independence.
\end{insight}

% =====================================================================
\section{The algorithmic local lemma}
% =====================================================================

The proof above is non-constructive: it certifies that a good outcome exists
without finding one among exponentially many possibilities. For decades this was
a genuine limitation, until a startlingly simple algorithm closed the gap.

\begin{theorem}{Moser--Tardos 2010}{moser-tardos}
In the variable model, if the general local-lemma condition holds, then the
following algorithm terminates in expected polynomial time and outputs an
assignment avoiding every $A_i$: while some $A_i$ is violated, pick one and
resample all the variables it depends on.
\end{theorem}

\begin{insight}
The rule ``resample whatever is currently broken'' looks too simple to
terminate, and the surprise is that it does. The analysis is an entropy-%
compression argument: a long run would let you reconstruct all the random bits
from a short log of which events were fixed, and that much compression is
impossible past a linear number of steps. So the existence the lemma guarantees
comes with an efficient algorithm --- the idea that made much of the probabilistic
method constructive.
\end{insight}

% =====================================================================
\section*{Exercises}
\addcontentsline{toc}{section}{Exercises}
% =====================================================================

\begin{enumerate}[itemsep=4pt]
  \item Show that the edges of $K_n$ can be coloured with at most $3\sqrt n$
        colours so that no triangle is monochromatic.
  \item Prove that every $k$-uniform $k$-regular hypergraph can be vertex-coloured
        with at most $k/\log k$ colours so that each colour appears $O(\log k)$
        times on every edge.
  \item $\bigstar$ (Hitting thin rectangles.) Prove there is $C>0$ so that for all
        small $\varepsilon$ one can pick one point in each unit grid square of
        $\mathbb R^2$ so that every $\varepsilon\times(C/\varepsilon)\log(1/\varepsilon)$
        rectangle (any orientation) contains a chosen point.
  \item (List colouring.) Prove there is $c>0$ so that if every vertex has a list
        of $k$ colours and each colour is acceptable to at most $ck$ neighbours of
        any vertex, then a proper list colouring exists.
  \item Prove that for every $\varepsilon>0$ there is a $\{0,1\}$ sequence
        $(a_i)$ and $\ell_0$ so that for all $\ell>\ell_0$ and all $i$, the blocks
        $(a_i,\dots,a_{i+\ell-1})$ and $(a_{i+\ell},\dots,a_{i+2\ell-1})$ differ in
        at least $(\tfrac12-\varepsilon)\ell$ coordinates.
  \item Prove that for every $\Delta$ there is $k$ so that any graph of maximum
        degree $\le\Delta$ has a $k$-colouring with no path
        $v_1v_2\cdots v_{2\ell}$ having $v_i$ the same colour as $v_{i+\ell}$ for
        all $i\le\ell$.
  \item Prove that every graph of maximum degree $\Delta$ can be properly
        edge-coloured so that no colour class contains a path of length $3$, using
        $O(\Delta)$ colours.
\end{enumerate}
